% ch9_c §9.6–9.7 (书页 313–323 / p327–p337)
%--B-TAIL--
%JOIN%
，其直径随 $H$ 的减小而增大。当这条轨道大到一定程度时，电子就会撞到薄片的下表面而被散射。表面阻抗对“有效”电子（§§ 8.7、9.2）颇为敏感，当式 (9.47) 中的 $2r_{x}$ 等于样品厚度时就会改变（图 166）。

这一现象完全没有利用回旋共振条件；而且由于表面是名副其实的空间不连续处，它比磁声效应更为明锐。但我们同样观察到在 $1/H$ 下具有恒定周期的振荡，来自各组极值轨道的贡献之间还发生干涉。这是因为，从薄片顶面趋肤深度出发的有效电子，会在顶面下方卡钳直径处汇合成一个薄薄的电流层。这个新的“趋肤层”又充当另一组轨道的源——如此等等。如果这样一条轨道链恰好能塞进薄片的厚度，从而在下表面产生一个“趋肤深度”区，样品对射频场就会显得比我们预料的透明得多。可是一旦下端轨道上的电子碰到这个表面，阻抗便急剧改变。再加上把不同的极值轨道组组合起来的种种可能性，观测结果变得相当复杂。不过，这些以及类似的非共振\emph{磁形态学}（magnetomorphic）现象，能够给出有关费米面的大量细致信息。

\section{轨道的量子化}

回旋频率 $\omega_{H}$ 就好比简谐振子的频率。因此我们预期会看到能量以 $\hbar\omega_{H}$ 为单位量子化。事实确实如此——尽管一个完全普遍的证明至今尚未给出。

但来考虑磁场中的自由电子。它们满足薛定谔方程
\begin{equation*}
\frac{1}{2m}\left(\frac{\hbar}{i}\nabla-\frac{e}{c}\mathbf{A}\right)^{2}\psi=\mathcal{E}\psi,
\tag{9.51}
\end{equation*}
其中 $\mathbf{A}$ 是矢势。我们要求出这个系统的定态。选取规范
\begin{equation*}
\mathbf{A}=(0,Hx,0),
\tag{9.52}
\end{equation*}
其旋度给出沿 $z$ 方向的场 $H$。

把式 (9.52) 代入式 (9.51)，我们必须求解
\begin{equation*}
\frac{\partial^{2}\psi}{\partial x^{2}}+\left(\frac{\partial}{\partial y}-\frac{ieH}{\hbar c}x\right)^{2}\psi+\frac{\partial^{2}\psi}{\partial z^{2}}+\frac{2m\mathcal{E}}{\hbar^{2}}\psi=0.
\tag{9.53}
\end{equation*}
这显然具有如下形式的解：
\begin{equation*}
\psi(x,y,z)=\exp\left\{i(\beta y+k_{z}z)\right\}u(x),
\tag{9.54}
\end{equation*}
其中 $u(x)$ 须满足方程
\begin{equation*}
\frac{\partial^{2}u}{\partial x^{2}}+\left\{\frac{2m\mathcal{E}'}{\hbar^{2}}-\left(\beta-\frac{eH}{\hbar c}x\right)^{2}\right\}u=0,
\tag{9.55}
\end{equation*}
其中
\begin{equation*}
\mathcal{E}'=\mathcal{E}-\frac{\hbar^{2}}{2m}k_{z}^{2}.
\tag{9.56}
\end{equation*}

沿 $z$ 方向——也就是平行于场的方向——的运动与自由电子完全相同，它对动能的贡献也一样。但对于 $(x,y)$ 平面内的运动，我们必须求解一个新的本征值方程 (9.55)，它可以写成
\begin{equation*}
-\frac{\hbar^{2}}{2m}\frac{\partial^{2}u(x)}{\partial x^{2}}+\frac{1}{2}m\left(\frac{eH}{mc}x-\frac{\hbar\beta}{m}\right)^{2}u(x)=\mathcal{E}'u(x).
\tag{9.57}
\end{equation*}
这不外乎就是简谐振子的波函数 $u(x)$ 的一维薛定谔方程，振子频率为
\begin{equation*}
\omega_{H}=\frac{eH}{mc},
\tag{9.58}
\end{equation*}
中心位于点
\begin{equation*}
x_{0}=\frac{1}{\omega_{H}}\frac{\hbar\beta}{m}.
\tag{9.59}
\end{equation*}
于是，
\begin{equation*}
\mathcal{E}'=\left(n+\frac{1}{2}\right)\hbar\omega_{H},
\tag{9.60}
\end{equation*}
而
\begin{equation*}
\mathcal{E}=\left(n+\frac{1}{2}\right)\hbar\omega_{H}+\frac{\hbar^{2}}{2m}k_{z}^{2}.
\tag{9.61}
\end{equation*}

这正是我们想要的结果：电子态的能量表为沿磁场方向的平移动能，加上垂直于磁场的平面内回旋运动的量子化能量之和。

有趣的问题是——我们怎样数状态？设有一个盒子，如 §1.7 中那样，边长为 $L_{x}$、$L_{y}$、$L_{z}$。显然，$k_{z}$ 像通常一样以 $2\pi/L_{z}$ 为单位量子化。同样，由式 (9.54) 的形式，$\beta$ 以 $2\pi/L_{y}$ 为单位量子化。但要注意，能量与 $\beta$ 无关，因此人们或许以为，对给定的 $n$，$\beta$ 可以取自一个无限集合中的任何值。

但情况并非如此。如式 (9.59) 所示，函数 $u$ 对 $\beta$ 的依赖体现在它以某点为中心：
\begin{equation*}
x_{0}=\frac{1}{\omega_{H}}\frac{\hbar\beta}{m}=\frac{v_{y}}{\omega_{H}}.
\tag{9.62}
\end{equation*}
这实际上意味着：如果电子以速度 $v_{y}$ 沿 $y$ 方向出发，它将在磁场中走一条圆形路径，圆心为 $x_{0}$。这条路径不能太大。$x_{0}$ 必须落在盒子之内，于是有
\begin{equation*}
0<x_{0}<L_{x}
\tag{9.63}
\end{equation*}
作为对这一中心允许位置的限制。

%FIG{167}: {磁场中电子薛定谔方程的解。}

但由式 (9.62)，这就限制了 $\beta$ 的允许取值范围。这个变量不仅以 $2\pi/L_{y}$ 为单位量子化，而且还必须满足
\begin{equation*}
0<\beta<\frac{m\omega_{H}}{\hbar}L_{x}=\frac{eH}{c\hbar}L_{x}.
\tag{9.64}
\end{equation*}
因此，只有
\begin{equation*}
p=\frac{L_{y}}{2\pi}\frac{m\omega_{H}}{\hbar}L_{x}
\tag{9.65}
\end{equation*}
个不同的 $\beta$ 值。式 (9.61) 的每个能级，对应于 $n$ 和 $k_{z}$ 的一种特定选择，是 $p$ 重简并的。

磁场当然已经破坏了我们原先的量子化图式。变量 $k_{x}$、$k_{y}$、$k_{z}$ 曾是我们在 k 空间中的坐标，如今不再是“好量子数”；波函数 (9.54) 也不再各自对应于这组参数的一组固定取值。尽管如此，让我们仍用这同一空间，并且把具有式 (9.61) 所给的某个 $\mathcal{E}$ 值的 $p$ 个新能级，用我们原图式中对应此能量的等能面来表示。略去 $z$ 坐标，这些面便是 $(k_{x},k_{y})$ 平面中的圆（译注：原文印作“$(x,y)$ 平面”。）。新的状态其实并不定域于圆上任何一点；它们将绕着圆以频率 $\omega_{H}$ 转动。但我们可以通过指明各能级所在的圆，来对磁场中的各个能级进行分类。

%FIG{168}: {自由电子的量子化图式：(a) 无磁场时；(b) 有磁场时。}

\emph{事实上我们可以核验：与 k 空间中任一给定宏观体积相联系的能级总数，在新图式中与原先相同。}我们可以利用式 (9.7)；相隔能量 $\delta\mathcal{E}$ 的两个等能面之间的面积为
\begin{equation*}
\delta\mathcal{A}=\frac{2\pi m_{H}^{*}}{\hbar^{2}}\delta\mathcal{E}.
\tag{9.66}
\end{equation*}
设 $\delta\mathcal{E}$ 是回旋频率的一个量子 $\hbar\omega_{H}$。回想在常规量子化图式中，“允许状态”的密度是每单位 $(k_{x},k_{y})$ 平面面积 $(L_{x}L_{y})/(2\pi)^{2}$（参见 §1.6）。于是，两条量子化轨道之间的状态数为
\begin{align*}
\frac{L_{x}L_{y}}{(2\pi)^{2}}\delta\mathcal{A}&=\frac{2\pi m_{H}^{*}}{\hbar^{2}}\hbar\omega_{H}\frac{L_{x}L_{y}}{(2\pi)^{2}}\\
&=p
\tag{9.67}
\end{align*}
即式 (9.65)。

这正是我们想要的结果。它表明，磁场的效应可以说是在 k 空间中造就出这些量子化轨道，并使自由电子状态“凝聚”到最近的一条这样的轨道上。每条轨道上的状态数，恰好等于它所处的那个环带内“允许状态”的数目。

那么，当我们处理的是晶体中的电子时，一般情形又会怎样？对单个能带，§6.5 的准经典理论有效。等效哈密顿量的薛定谔方程 (6.30) 在磁场中的解，在大量子数下将满足对应原理极限，并将按 Bohr 相位积分公式量子化：
\begin{equation*}
\oint\mathbf{p}\cdot d\mathbf{r}=(n+\gamma)\,2\pi\hbar,
\tag{9.68}
\end{equation*}
其中 $n$ 是整数，$\gamma$ 是相位修正（典型值为 $\frac{1}{2}$），$\mathbf{p}$ 和 $\mathbf{r}$ 是共轭变量，表示粒子描出其轨道时的动量和位置。

我们可以利用式 (9.45) 来计算 $\mathbf{r}$，即电子在真实空间中沿其路径的位置；并且可以采用通常的动量规则
\begin{equation*}
\mathbf{p}=\hbar\mathbf{k}+\frac{e}{c}\mathbf{A}.
\tag{9.69}
\end{equation*}
由简单的几何推理即得：\emph{轨道在 k 空间中的面积是量子化的}：
\begin{equation*}
\mathcal{A}_{n}=\frac{2\pi eH}{c\hbar}(n+\gamma).
\tag{9.70}
\end{equation*}
由式 (9.61)、(9.66) 和 (9.5) 显然可见，这正是前面已对自由电子得到的结果。

要在磁场中描述电子能级，现在需要在 k 空间中作如下构造。选定一个 $n$ 值。在费米面每个垂直于磁场的截面上，画出面积为 $\mathcal{A}_{n}$ 的等能围线。把这些围线连成一根连续的管子，其轴平行于 $H$，且横截面积处处相同。对其余的 $n$ 值同样画出管子。假定常规量子化图式中的“允许点”已全部凝聚到最近的管子上。这就给出了在能量上简并的状态数，它们现在应被想象成以各自轨道相应的回旋频率绕每根管子环流。

%FIG{169}: {量子化磁能级的管。}

\section{de Haas--van Alphen 效应}

这种量子化有一些有趣而惊人的效应。电子气体作为一个整体的能量依赖于磁场的强度。而磁化率——它基本上是抗磁性的（这里我们忽略自旋效应），并且与温度无关——却随着磁场的改变而振荡。这就是 de Haas--van Alphen 效应。

要对一般形状的费米面计算这个效应，我们可以按如下步骤进行。先写下系统的自由能；对 Fermi--Dirac 系集（assembly），它由下式给出：
\begin{equation*}
F=N\zeta-kT\sum_{i}\ln\left(1+e^{(\zeta-\mathcal{E}_{i})/kT}\right),
\tag{9.71}
\end{equation*}
其中对所有可能的状态求和，状态能量为 $\mathcal{E}_{i}$；但我们由被占据状态的数目来固定 $\zeta$，即 Fermi 势（例如如式 (4.9) 中那样）。

对我们的量子化 \emph{Landau 能级}系统，能量由下式给出：
\begin{equation*}
\mathcal{E}_{i}=\mathcal{E}(n+\gamma,k_{z}),
\tag{9.72}
\end{equation*}
其中 $n$ 是轨道量子数，如式 (9.61) 中那样——但当然，如式 (9.65) 那样，每个这样的能级是 $p$ 重简并的。取一块单位边长的晶体立方；按式 (9.67) 或 (9.70)，k 空间中每根管子长度为 $dk_{z}$（照常沿 $H$ 方向量度）的一段上将有
\begin{equation*}
\frac{1}{4\pi^{3}}d^{3}k=\frac{eH}{2\pi^{2}\hbar c}\,dk_{z}
\tag{9.73}
\end{equation*}
个状态属于它。于是
\begin{equation*}
F=N\zeta-kT\int_{-\infty}^{\infty}\left\{\frac{eH}{2\pi^{2}\hbar c}\sum_{n=0}^{\infty}\ln\left(1+\exp\left[\left\{\zeta-\mathcal{E}(n+\gamma,k_{z})\right\}/kT\right]\right)\right\}dk_{z}.
\tag{9.74}
\end{equation*}

要算这个讨厌的积分，我们用一招纯数学的技巧——Poisson 求和公式。考虑 $f(x)$，一个任意函数。在 $n<x<n+1$ 范围内可以写成
\begin{equation*}
f(x)=\sum_{s=-\infty}^{\infty}e^{-2\pi ixs}g_{s},
\tag{9.75}
\end{equation*}
其中
\begin{equation*}
g_{s}=\int_{n}^{n+1}f(x)\,e^{2\pi ixs}dx.
\tag{9.76}
\end{equation*}
这正是一个 Fourier 级数，如式 (1.7)。于是可以写
\begin{align*}
f\left(n+\frac{1}{2}\right)&=\sum_{s=-\infty}^{\infty}e^{-2\pi is(n+\frac{1}{2})}g_{s}\\
&=\sum_{s=-\infty}^{\infty}e^{-i\pi s}g_{s}\\
&=\sum_{s=-\infty}^{\infty}(-1)^{s}\int_{n}^{n+1}f(x)\,e^{2\pi ixs}dx.
\tag{9.77}
\end{align*}
现在把上式对变量 $x$ 的所有区间求和，得
\begin{equation*}
\sum_{n=0}^{\infty}f\left(n+\frac{1}{2}\right)=\int_{0}^{\infty}f(x)\,dx+2\sum_{s=1}^{\infty}(-1)^{s}\int_{0}^{\infty}f(x)\cos 2\pi xs\,dx.
\tag{9.78}
\end{equation*}
我们把这个公式用到式 (9.74) 中对 $n$ 的求和上：
\begin{align*}
F&=N\zeta-kT\int_{-\infty}^{\infty}\left\{\frac{eH}{2\pi^{2}\hbar c}\int_{0}^{\infty}\ln\left(1+\exp\left[\left\{\zeta-\mathcal{E}(x+\gamma,k_{z})\right\}/kT\right]\right)dx\right\}dk_{z}\\
&\quad-2kT\sum_{s=1}^{\infty}\int_{-\infty}^{\infty}\left\{\frac{eH}{2\pi^{2}\hbar c}(-1)^{s}\int_{0}^{\infty}\ln\left(1+\exp\left[\left\{\zeta-\mathcal{E}(x+\gamma,k_{z})\right\}/kT\right]\right)\right.\\
&\qquad\left.\times\cos 2\pi xs\,dx\right\}dk_{z}.
\tag{9.79}
\end{align*}

我们暂且把注意力集中于如下形式的积分：
\begin{align*}
&\int_{0}^{\infty}\ln\left(1+\exp\left[\left\{\zeta-\mathcal{E}(x+\gamma,k_{z})\right\}/kT\right]\right)\cos 2\pi xs\,dx\\
&\quad=\frac{1}{4\pi^{2}s^{2}kT}\left[f^{0}(\mathcal{E})\frac{\partial\mathcal{E}}{\partial x}\right]_{x=0}+\int_{0}^{\infty}\frac{\cos 2\pi xs}{4\pi^{2}s^{2}kT}\left[f^{0}(\mathcal{E})\frac{\partial^{2}\mathcal{E}}{\partial x^{2}}+\frac{\partial f^{0}}{\partial x}\frac{\partial\mathcal{E}}{\partial x}\right]dx.
\tag{9.80}
\end{align*}
这里我们用到了
\begin{equation*}
kT\,\frac{\partial\ln\left(1+e^{(\zeta-\mathcal{E})/kT}\right)}{\partial\mathcal{E}}=f^{0}(\mathcal{E}),
\tag{9.81}
\end{equation*}
即式 (4.8) 所定义的 Fermi 函数，并且已分部积分两次。我们将略去式 (9.80) 中的第一项，因为我们只关心自由能的振荡部分。利用式 (9.61)，按定义有
\begin{equation*}
\mathcal{E}(x+\gamma,k_{z})=(x+\gamma)\hbar\omega_{H}+f(k_{z}),
\tag{9.82}
\end{equation*}
于是 $\partial\mathcal{E}/\partial x=\hbar\omega_{H}$，而 $\partial^{2}\mathcal{E}/\partial x^{2}=0$。再者，积分下限不妨取 $-\infty$，因为对积分的重要贡献只来自费米能级 $\mathcal{E}=\zeta$ 附近。于是，费米分布的这一薄层对自由能的贡献正比于
\begin{equation*}
I_{s}(k_{z})=\frac{\hbar\omega_{H}(\zeta,k_{z})}{4\pi^{2}s^{2}kT}\int_{-\infty}^{\infty}\cos 2\pi xs\,\frac{\partial f^{0}}{\partial x}\,dx.
\tag{9.83}
\end{equation*}

这个积分虽然名义上具有式 (4.13) 的形式，却不能用那个公式求值，因为被积函数在 Fermi 层厚度之内振荡得太快。但我们可以利用这样一个事实：$-\partial f^{0}/\partial x$ 在点 $X$ 处取极大，$X$ 满足
\begin{equation*}
\mathcal{E}(X,k_{z})=\zeta:
\tag{9.84}
\end{equation*}
我们可以说
\begin{equation*}
X=\frac{c\hbar}{2\pi eH}\,\mathcal{A}(\zeta,k_{z}),
\tag{9.85}
\end{equation*}
其中 $\mathcal{A}(\zeta,k_{z})$ 是费米面在切片 $k_{z}$ 处、实际费米能级上的横截面面积。我们把积分在这个点附近展开——记住 $\partial f^{0}/\partial x$ 是 $x$ 的偶函数——便得
\begin{align*}
I_{s}(k_{z})&\approx\frac{\hbar\omega_{H}}{4\pi^{2}s^{2}kT}\int_{-\infty}^{\infty}\cos 2\pi s\left(X+\frac{kT}{\hbar\omega_{H}}\eta\right)\frac{\partial f^{0}}{\partial\eta}\,d\eta\\
&=\frac{\hbar\omega_{H}}{4\pi^{2}s^{2}kT}\cos 2\pi sX\int_{-\infty}^{\infty}\cos\left(\frac{2\pi skT}{\hbar\omega_{H}}\eta\right)\frac{\partial f^{0}}{\partial\eta}\,d\eta\\
&=-\frac{\hbar\omega_{H}}{4\pi^{2}s^{2}kT}\cos\left(\frac{sc\hbar\mathcal{A}(\zeta,k_{z})}{eH}\right)\left\{\frac{2\pi^{2}skT/\hbar\omega_{H}}{\sinh\left(2\pi^{2}skT/\hbar\omega_{H}\right)}\right\}\\
&=-g(s,k_{z})\cos\left(\frac{sc\hbar\mathcal{A}(\zeta,k_{z})}{eH}\right),\quad\text{姑且记之。}
\tag{9.86}
\end{align*}

写到这里，我们或许应该停下来喘口气，看看究竟发生了什么。我们考虑的是费米面的一个特定切片。由整个图形的体积所固定的费米能级，未必恰好与某条量子化磁轨道重合。随着场的改变，量子化轨道被拉进、拉出，扫过费米能级，如图 170 所示。

%FIG{170}: {(a) 费米面未必与量子化轨道重合。(b) 磁场改变时，量子化能级扫过费米能级。}

%FIG{171}: {磁场改变时磁能级的占据情况。}

设我们处于 $T=0$，费米分布完全锐利。当一条量子化轨道扫过 $\zeta$ 时会发生什么？例如，设场的大小恰使 $\zeta$ 位于两条轨道的正中间（图 171(a)）。这时费米能级以下的状态数，与完全没有磁能级时的情形一模一样——但电子气体的总能量将低于没有磁场时的值——费米能级上每个电子约低 $\frac{1}{2}\hbar\omega_{H}$ 的量级。现在，随着 $H$ 增大，这些电子将被拉向费米能级（图 171(b)），其自由能随之升到一个极大。而当磁能级穿过费米能级时，它开始倒空（图 171(c)），电子的平均能量再度下降，并在 $\zeta$ 再次位于两条量子化轨道正中间时达到极小。这样，电子气体的自由能便规则地振荡，其周期由量子化轨道与费米能级相邻两次重合的间隔决定。这就是式 (9.86) 的含义；周期来自条件
\begin{equation*}
\frac{sc\hbar\mathcal{A}(\zeta,k_{z})}{eH}=2n\pi.
\tag{9.87}
\end{equation*}

从式 (9.74) 到式 (9.81) 的大部分代数演算，都是为把振荡按 Fourier 级数来分析而设的手段，因为每个能级扫过费米面时能量的变化并不是简单的正弦函数。这也解释了指标 $s$ 的由来——它给出更高的谐波。

我们也已经考虑到 $T\neq 0$ 时费米面并非完全锐利这一事实；事实上，若
\begin{equation*}
kT\gg\hbar\omega_{H},
\tag{9.88}
\end{equation*}
则所有振荡都将被抹平。换句话说，磁场必须强到使轨道的间距达到费米面上热层厚度的量级。

但分析到此还没有完结。以上只是对 $k_{z}$ 处一个特定截面的计算；我们必须对所有不同的切片求和：
\begin{equation*}
F=2kT\sum_{s=1}^{\infty}(-1)^{s}\int_{-\infty}^{\infty}\frac{eH}{2\pi^{2}\hbar c}\,g(s,k_{z})\cos\left(\frac{sc\hbar\mathcal{A}(\zeta,k_{z})}{eH}\right)dk_{z}.
\tag{9.89}
\end{equation*}
这是一个 Fresnel 型积分；众所周知，主要贡献来自相位驻定的区域。设
\begin{equation*}
\frac{\partial\mathcal{A}(\zeta,k_{z})}{\partial k_{z}}=0
\tag{9.90}
\end{equation*}
在 $k_{z}=k_{0}$ 处成立。在该点附近展开：$k_{z}=k_{0}+k'$，且
\begin{equation*}
\mathcal{A}=\mathcal{A}_{0}\pm\frac{1}{2}k'^{2}\mathcal{A}_{0}''\dots
\tag{9.91}
\end{equation*}
（$g(s,k_{z})$ 的变化可以忽略——而且我们也不言自明地假定了费米面的每一部分都有一个对称中心。）于是
\begin{align*}
F&=2kT\sum_{s=1}^{\infty}(-1)^{s}\frac{eH}{2\pi^{2}\hbar c}\,g(s,k_{0})\int_{-\infty}^{\infty}\cos\left\{\frac{sc\hbar}{eH}\left(\mathcal{A}_{0}\pm\frac{1}{2}k'^{2}\mathcal{A}_{0}''\dots\right)\right\}dk'\\
&\approx 2kT\sum_{s=1}^{\infty}(-1)^{s}\left(\frac{eH}{2\pi^{2}\hbar c}\right)g(s,k_{0})\left(\frac{2eH}{sc\hbar|\mathcal{A}_{0}''|}\right)^{\frac{1}{2}}\cos\left(\frac{sc\hbar\mathcal{A}_{0}}{eH}\pm\frac{\pi}{4}\right)\\
&=2kT\sum_{s=1}^{\infty}(-1)^{s}\left(\frac{eH}{2\pi sc\hbar}\right)^{\frac{1}{2}}\frac{1}{\sinh\left(2\pi^{2}skT/\hbar\omega_{H}\right)}\,|\mathcal{A}_{0}''|^{-\frac{1}{2}}\cos\left(\frac{sc\hbar}{eH}\mathcal{A}_{0}\pm\frac{\pi}{4}\right),
\tag{9.92}
\end{align*}
这里利用了 Fresnel 积分的标准性质。对 $H$ 微商两次，就可以把它表示为对金属磁化率的一个贡献。

这个公式看上去非常复杂——但其解释同样十分简单。我们不去盯住单个截面，而是注视整根磁管随磁场增大而向外生长的全过程。如上面所见，每根管穿过费米能级时能量会发生变化。但当一根给定的管子生长时，所发生的不过是它与费米面的交线沿管子上下移动（见图 169），因此对能量几乎没有影响——除非管子整个脱离费米面，这发生在横截面积取极大或极小之处。这解释了 $\mathcal{A}_{0}$ 出现在周期因子中的缘由。对一次谐波，
\begin{equation*}
\frac{c\hbar}{eH}\mathcal{A}_{0}=2\pi(n+\gamma),
\tag{9.93}
\end{equation*}
其中 $\gamma$ 是相位修正。于是，\emph{把磁矩作为 $1/H$ 的函数作图时，振荡的周期就直接给出费米面垂直于磁场方向的最大或最小截面的面积。}

实际上，对于复杂的费米面，如果同时存在几个不同的截面极大与极小，每一处都会贡献自己的振荡项，合成的磁矩作为场的函数就可能表现出非常复杂的行为。每个振荡的振幅将依赖于 $\mathcal{A}_{0}''$——极值截面附近费米面的局域曲率。它还依赖于温度，近似为 $\exp(-kT/\hbar\omega_{H})$，由此可以估算驻定轨道上的回旋频率。但如果存在杂质散射，振幅还会被再乘上一个形如
